%!TEX root = TheoretischePhysikC.tex


Information über die Organisation, die \emph{Studienleistung} und \emph{Prüfungsleistung}: Siehe unbedingt Stud.ip.

Wir müssen ein ganzes Spektrum von Themen abdecken und zwar
\begin{itemize}
\item Experimentelle Befunde
\item Postulate der Quantenmechanik
\item Unschärfe
\item Schrödingergleichung
\item einfache Potentialprobleme
\item harmonischer Oszillator
\item Wasserstoff-Atom
\item identische Teilchen
\item Verschränkung
\item Hauptsätze der Thermodynamik
\item Carnot'sche Prozesse
\item Statistische Beschreibung des Gleichgewichts
\item Ideale (Quanten-)Gase
\end{itemize}

\emph{,,Bug bounty``}: Ich hätte es gern, dass das Skript so fehlerfrei wie möglich ist. Deshalb habe ich ein ,,Bug bounty programme`` eingeführt. Wenn Sie einen Fehler entdecken, schreiben Sie mir bitte eine Email mit den Einzelheiten. Ich werde alle solche Beiträge aufnehmen und am Ende des Semesters drei Preise vergeben. Für den ersten Platz bekommt man eine Flasche guten australischen Weins (beziehungsweise ein vergleichbares alkoholfreies Getränk). Für den zweiten Platz bekommt man zwei Packungen australischer Kekse und für den dritten Platz eine Packung Kekse. Beispiele für ,,Bugs`` sind Vorzeichen-Fehler und grammatikalische Fehler/Unklarheiten.

Wir fangen mit ein paar historischen Bemerkungen an. 

\section{Vorlesung: Der Geburt der Quantenmechanik}
Wir befinden uns in einem Quanten-Universum. Allerdings haben wir bis 1900 die Quantenaspekte des Universums nicht bemerkt, weil uns die klassische Physik eine sehr gute Annäherung der Physik des Alltags geliefert hat. Die Quantenmechanik wäre viel einfacher zu verstehen, wenn wir Quantenaspekte jeden Tag beobachten könnten! (Eine der größten Herausforderungen beim Erlernen der Quantenmechanik, ist zu verstehen, warum die Näherung durch die klassische Physik so gut funktioniert.)

Vor kurzem konnten wir eine sensationelle Entwicklung in der Physik beobachten: Die ersten Quantencomputer sind entwickelt worden. Jetzt haben wir die Möglichkeit, zum ersten Mal, direkt mit Quantensystemen zu experimentieren (mit Hilfe von Cloud Services, z.B. IBM Q experience). Quantenaspekte werden langsam sichtbar werden. 

Bevor wir uns mit den neuesten Ergebnissen in der Anwendung befassen können, betrachten wir zunächst die Entwicklung der QM und ihren formalen Aufbau. 

\subsection{Vorgeschichte}
Die Vorgeschichte der Quantenmechanik umfasst eine Reihe von Beobachtungen, die die These, dass unser Universum durch die klassische Mechanik beschrieben wird, allmählich in Frage stellten. Einige der wichtigsten dieser Beiträge sind
\begin{enumerate}
	\item \emph{Schwarzkörperstrahlung}. Das Spektrum der emittierten Strahlung ist nicht mit klassischer Physik erklärbar. Die Annahme von Lichtquanten mit Energie $E=h\nu$ liefert das richtiges Spektrum. Die konstante ,,h`` (ursprünglich ,,Hilfsgröße``) wurde von Planck zur Herleitung des Strahlungsgesetzes:
	\begin{equation*}
		\text{Energiedichte der Strahlung im Intervall $d\nu$} = U(\nu)d\nu = \frac{8\pi h \nu^3}{c^3}\frac{d\nu}{e^{\frac{h\nu}{kT}}-1}
	\end{equation*}
	postuliert (1899). Heute wird in der theoretischen Physik vorwiegend die \emph{reduzierte Planck-Konstante} 
	\begin{equation*}
		\eqbox{\hbar = \frac{h}{2\pi}}
	\end{equation*}
	verwendet.
	\item \emph{Photoeffekt}. Die Energie der durch Licht ausgelösten Elektronen (Photoelektronen) ist nur durch die Lichtfrequenz bestimmt (Lenard, 1899; Einstein 1905):
	\begin{equation*}
		\eqbox{E_{\text{el}} = \frac12 m\nu^2 = h\nu -W}
	\end{equation*}
	mit Bindungsenergie $W$. Diese Beobachung war ein Hinweis auf \emph{Lichtquanten} mit Energie $E=h\nu$.
	\item \emph{Stabilität der Atome}, \emph{Linienspektren}, \emph{Franck-Hertz-Versuch} ($\sim$1911). In der klassischen Physik wären Elektronenbahnen in einem Atom wegen Strahlungsverlusten instabil. Strahlung wäre kontinuierlich über alle Frequenzen verteilt. Linienspektren und Franck-Hertz-Versuche deuten darauf hin, dass in Atomen nur bestimmte, diskrete Energieniveaus erlaubt sind. 
	\item \emph{Compton-Effekt}. Die Streuung von Licht an Elektronen wird durch Annahme von Lichtquanten mit $E=h\nu$ und Impuls $p=h/\lambda$ sowie Energie- und Impulserhaltung richtig beschrieben. 
	\item \emph{Stern-Gerlach-Versuch} (1922). Der Drehimpuls hat endlich viele Möglichkeiten der Ausrichtung relativ zu einem angelegten Magnetfeld (Richtungsquantelung).
	\item \emph{Interferenz bei Elektronen} (Davisson-Germer-Experiment, 1929; Beugung an Kristallen). Dieses Experiment lieferte einen Hinweis auf die Welleneigenschaften des Elektrons (postuliert 1924 von de Broglie).
\end{enumerate}

\begin{center}
\parbox[c][5cm][c]{14cm}{\begin{tcolorbox}[colback=banana,arc=0cm,frame hidden,enhanced]Materie und Strahlung besitzen sowohl Wellen- als auch Teilcheneigenschaften.\end{tcolorbox}}
\end{center}



Das Geburtsjahr der Quantenmechanik ist 1926. In diesem Jahr veröffentlichte Schrödinger eine spannende Reihe von Arbeiten über die Wellenmechanik und Quantisierung. Ein paralleler Zugang begann 1925 mit Überlegungen von Heisenberg über die Messbarkeit von Phasenbeziehungen von Teilchen im Bohrschen Atommodell. Diese Arbeit hat zwei Teams zur Aufstellung von Versionen der Quantenmechanik inspiriert. Einerseits Born und Jordan mit Heisenberg zur sogenannten Matrizenmechanik und andererseits Dirac. Noch 1926 zeigt Schrödinger die Äquivalenz beider Zugänge. Die statistische Interpretation der Theorie wird von Born (ebenfalls 1926) geliefert. Damit ist die Grundlage der Theorie abgeschlossen und Begriffe wie ,,Welle-Teilchen Dualismus`` sind entstanden. 

\begin{center}
\parbox[c][5cm][c]{14cm}{\begin{tcolorbox}[colback=banana,arc=0cm,frame hidden,enhanced]Die Quantenmechanik ist eine \emph{probabilistische Theorie}. Sie liefert uns somit nur Vorhersagen in Form von \emph{Wahrscheinlichkeiten}.\end{tcolorbox}}
\end{center}

Fast gleichzeitig wurde die Theorie in ihre heutige mathematische Form gebracht. So ließ sich Born von Norbert Wiener dabei beraten, aus Matrizen Operatoren zu machen und so auch ,,aperiodische Probleme`` ohne vorgegebene Basis (z.B.\ ein freies Teilchen) zu behandeln. Hilbert hielt im WS 1926/27 in Göttingen eine Vorlesung über Quantenmechanik. Nebenbei führte sein junger Assistent John von Neumann den Hilbertraum, Dichtematrizen und die Entropie für Quantenzustände mit einer Publikation in den Göttinger Mathematischen Nachrichten ein. Diese wurde von Physikern zunächst nicht gelesen. Die Rolle der Gruppentheorie wurde von Wigner und Weyl geklärt.

Bemerkenswert ist, dass trotz der phantastischen Erfolge der Theorie die Interpretation umstritten blieb. Trotz vielen Ideen und deren Entwicklungen, ist die Interpretation der QM immer noch kontrovers. Dennoch ist die QM heute allgemein akzeptiert und liefert uns das beste physikalische Modell, das wir in der Physik haben; es gibt bis heute keinen Widerspruch zur QM. 

Wir verlassen nun die spannende Geschichte der Quantenmechanik. Unsere Annahme ist, dass die Quantenmechanik die korrekte Beschreibung des Universums ist.

Unser nächstes Ziel ist es, die mathematischen Grundlagen der Quantenmechanik zusammenzufassen. Um dieses Ziel zu erreichen, müssen wir erst die \emph{Wahrscheinlichkeit} mathematisch behandeln. 

\subsection{Wahrscheinlichkeit}
Die Zahlen, welche die Quantenmechanik als Ergebnis liefert, sind \emph{Wahrscheinlichkeiten}. Wie sind solche Größen zu verstehen und wie arbeiten wir damit? Die \emph{Wahrscheinlichkeitstheorie} ist der richtige Rahmen, um zufällige Ereignisse und ihre Interpretation zu verstehen. Diese Theorie wird demnächst aufgebaut. 


\subsubsection{Stichprobenräume, Maße, Momente}
In der Wahrscheinlichkeitstheorie wird jedem System eine Menge $X$ zugeordnet, die man \emph{Elementarereignisraum} oder \emph{Stichprobenraum} nennt. 

\begin{beispiel}
	Nehmen wir den einfachsten Fall, in dem $X$ eine endliche Menge ist.Menge ist. Beispielsweise bei einem Würfel besteht der Elementarereignisraum aus den möglichen Augenzahlen. Somit gilt $X = \{1, 2, 3, 4, 5, 6\}$.
\end{beispiel}

\begin{übung}
	Beschreiben Sie die Menge $X$, die zu dem Elementarereignisraum \emph{zweier} Würfeln zugeordnet sein soll.  
\end{übung}

Eine \emph{Präparation} wird durch die Angabe einer Wahrscheinlichkeit $p(x)$ für jedes $x\in X$ charakterisiert. 

\begin{beispiel}
	Für einen fairen Würfel ist $p(x) = 1/6$ für alle $x$. Aber es gibt auch andere Wahrscheinlichkeiten, welche dann die Präparationen von \emph{gezinkten Würfeln} beschreiben.
\end{beispiel}

Die einzigen Einschränkungen an die Wahrscheinlichkeitsverteilung $p(x)$ sind
\begin{equation}
	p(x) \ge 0 \quad \text{und}\quad \sum_{x\in X} p(x)  = 1.
\end{equation}

Eine \emph{Messung} (Observable) wird durch eine Funktion 
\begin{equation}
	f:X\rightarrow \mathbb{R}
\end{equation}
beschrieben. 
\begin{beispiel}
	Für einen Würfel ist die \emph{Augenzahl} $f(x) = x$ eine Messung.
\end{beispiel}
Solche Messungen werden in der Wahrscheinlichkeitstheorie als \emph{Zufallsvariablen} bezeichnet. Der Wert $f(x)$ wird als der Messwert interpretiert, der sich ergibt, wenn man sicher weiß, dass das Elementarereignis $x$ vorliegt. Wenn dieses Ereignis bei $N$ Messungen vorliegt ($Np(x)$), ergibt sich die Summer der Messwerte mit $\sum_x N p(x) \cdot f(x)$. Dividiert durch N erhalten wir die Formel für den Erwartungswert:
\begin{equation}
	\text{Erwartungswert} = \mathbb{E}[f] = \sum_{x\in X} p(x)f(x).
\end{equation}

Eine wichtige Rolle spielt in der Quantenmechanik der Spezialfall der Zufallsvariablen $f$, die nur Ergebnisse $0$ oder $1$ liefern können. Solche Messungen dienen als Modell für Experimente, z.B. Photonendetektoren, die das Vorhandensein oder Nichtvorhandensein eines Teilchens messen. Diese Messungen sind durch eine Menge $S=\{x\in X\,|\, f(x) =1\}$ charakterisiert. Wir verwenden oft die Notation 
\begin{equation}
	\chi_S(x) = \begin{cases}
		1, \quad x\in S, \\
		0, \quad \text{sonst.}
	\end{cases}
\end{equation}
Solche Funktionen nennt man \emph{Indikatorfunktionen}. Die Teilmengen $S\subset X$
nennt man auch Ereignisse (nicht mehr Elementarereignisse). Der Erwartungswert 
\begin{equation}\label{eq:indicatorfunction}
	\mathbb{E}[\chi_S] = \sum_{x\in X} p(x)\chi_S(x) = \sum_{x\in S} p(x) \equiv \mathbb{P}[S]
\end{equation}
stellt die Wahrscheinlichkeit dar, dass das Ereignis $x\in S$ stattfindet.

Die Funktion $\mathbb{P}$ hat charakteristische Eigenschaften, die sich sofort aus der Definition ergeben. So gilt für den Fall dass die Teilmengen $S_j\subset X$ disjunkt sind, dass 
\begin{equation}
	\mathbb{P}\left[\bigcup_j S_j \right] = \sum_{j} \mathbb{P}[S_j].
\end{equation}
Mathematisch gesprochen: $P$ ist ein Maß, bzw. ein Wahrscheinlichkeitsmaß, wenn man die Normierung $\mathbb{P}[X] = 1$ betonen will. Wenn man die \emph{Mengenfunktion} $\mathbb{P}$ kennt, kann man auch das Erwartungswertfunktional $\mathbb{E}$ ausrechnen. 

Im Fall, dass $X$ eine endliche Menge ist, können wir das Erwartungswertfunktional wie folgt berechnen
\begin{equation}
	\mathbb{E}[f] = \sum_{x\in X} \mathbb{P}[\{x\}] f(x).
\end{equation} 
Das heißt, $p(x) = \mathbb{P}[\{x\}]$. 
\begin{übung}
	Sei $X = \{1,2,3\}$ und $p(x) = 1/3$ $\forall x\in X$. Berechnen Sie die Wahrscheinlichkeit $\mathbb{P}[S]$, im Fall $S=\{1,2\}$ ist. Sei $f(x) = x^2$. Wie lautet der Erwartungswert
	\begin{equation}
		\mathbb{E}[x^2] = ?
	\end{equation}
\end{übung} 

Wenn $X = \mathbb{R}$, verwenden wir eine andere Notation:
\begin{equation}
	\mathbb{E}[f] = \int \mathbb{P}(dx)\, f(x).
\end{equation}
(Die Summe wird durch ein Integral ersetzt.)
Diese Notation wird oft auch als 
\begin{equation}\label{eq:ctsprob}
	\mathbb{E}[f] = \int p(x) f(x) \, dx
\end{equation}
geschrieben, wobei $p(x)\,dx \equiv \mathbb{P}(dx)$ die Wahrscheinlichkeitsverteilung darstellt.

\begin{beispiel}
	Die Wahrscheinlichkeitsverteilung
	\begin{equation}
		p(x) \equiv \frac{1}{\sqrt{2\pi}}e^{-\frac{1}{2}x^2}
	\end{equation}
	stellt eine \emph{Normal-} oder \emph{Gauß-Verteilung} dar. Die Wahrscheinlichkeit, dass $x\in S$, wobei $S= \{x\,|\, x\ge 0\}$, wird durch das folgende Integral berechnet (nach Gl.~\ref{eq:indicatorfunction} und Gl.~\ref{eq:ctsprob}):
	\begin{equation}
		\begin{split}
			\mathbb{P}[S] = \mathbb{E}[\chi_S] &= \frac{1}{\sqrt{2\pi}}\int_{-\infty}^\infty \chi_S(x) e^{-\frac{1}{2}x^2}\, dx \\ &= \frac{1}{\sqrt{2\pi}}\int_{-\infty}^0 \chi_S(x) e^{-\frac{1}{2}x^2}\, dx + \frac{1}{\sqrt{2\pi}}\int_{0}^\infty \chi_S(x) e^{-\frac{1}{2}x^2}\, dx = \frac{1}{2}.
		\end{split}
	\end{equation}
\end{beispiel}

Für Maße auf $\mathbb{R}$ ist es oft besonders einfach, die Erwartungswerte der Koordinatenfunktion und ihrer Potenzen zu berechnen. Diese werden \emph{Momente} genannt. Für ein Wahrscheinlichkeitsmaß $\mathbb{P}(dx) = p(x)dx$ auf $\mathbb{R}$ nennt man
\begin{equation}
	M_n = \int x^n p(x)\, dx 
\end{equation}
das \emph{$n$-te Moment} von $\mathbb{P}$. Es ist immer $M_0 = 1$, wegen der Normierung auf Gesamtwahrscheinlichkeit 1. (Warnung: Alle anderen Momente brauchen nicht zu existieren, bzw.\ ergeben unter Umständen divergente Integrale.) 

Die praktisch wichtigsten Momente sind $M_1$ (\emph{Mittelwert}), und $M_2$. In der Praxis wird das zweite Moment oft mit der \emph{Varianz} ersetzt:
\begin{equation}
	\text{Var}[\mathbb{P}] \equiv \mathbb{E}[(x-M_1)^2] \equiv \int (x^2-2M_1 x +M_1^2) p(x)\, dx \underset{\text{Übung!}}{=} M_2-M_1^2.
\end{equation}
Aufgrund der Definition ist klar, dass Varianzen stets $\ge 0$ sind und nur dann verschwinden, wenn $\mathbb{P}$ auf einen Punkt konzentriert ist.